\chapter{Other word classes}
\label{chap:modifiers}


\section{Adjectives}

\subsection{Declension}

Adjectives must agree in gender and number with the noun they modify. Each adjective thus has four main forms: masculine singular, feminine singular, neuter singular, and plural. The plural form is used for all plural nouns, regardless of gender. Note that adjectives are not marked for case. Most adjectives follow the declension pattern of the adjective \nikp{bang}{good}, shown in Table~\ref{tab:adj-declension} below.

Agreement endings preserve lexical stress. The forms \nk{bang}, \nk{banger}, \nk{bange}, and \nk{banga} all have initial stress (see Chapter~\ref{chap:phonology}, \S\ref{sec:stress-prosody}).

Ordinary productive agreement also preserves the base's accent category unless an inherited paradigm specifies an alternation. Its surface contrast depends on the licensing of the primary-stressed rhyme. An unstressed agreement ending has no independent lexical accent and does not supply a missing tonal position for the stem (Chapter~\ref{chap:phonology}, \S\ref{sec:accent-licensing}).

\begin{table}[!ht]
    \caption{Declension of the adjective \nikp{bang}{good}}
    \label{tab:adj-declension}
    \small
    \begin{tabularx}{\linewidth}{XXXXX}
        \toprule
        \rowcolor{grammarTableHead}
         & \textbf{Masculine} & \textbf{Feminine} & \textbf{Neuter} & \textbf{Plural} \\
        \midrule
        \nk{bang} & \nk{bang} & \nk{banger} & \nk{bange} & \nk{banga}\\
        \bottomrule
    \end{tabularx}
\end{table}

The feminine ending \nk{-er} and the plural ending \nk{-a} are both pronounced [\ipa{ɐ}] in ordinary speech. Thus \nk{banger} and \nk{banga} are homophonous, while retaining their distinct agreement functions. The neuter ending \nk{-e} is realised through a syllabic nasal in \nk{bange}, giving [\ipa{ŋ̩}] at the end of the word; careful speech retains [\ipa{ə}]. The masculine \nk{bang} has no additional syllabic interval. With obstruent-final stems, neuter \nk{-e} undergoes apocope, subject to the cluster restrictions in Chapter~\ref{chap:phonology}, \S\ref{sec:apocope}.

Agreement is determined by the noun's gender and number even when two adjective forms have the same pronunciation. Articles and other elements of the noun phrase provide additional gender and number marking. The written four-form paradigm remains the same.

Adjectival participles are declined as follows:

\begin{table}[!ht]
    \caption{Declension of adjectival participles}
    \label{tab:part-declension}
    \small
    \begin{tblr}{
        colspec={XXXXX}
    }
        & Masculine & Feminine & Neuter & Plural \\
        Present & \nk{hjaltade} & \nk{hjaltadär} & \nk{hjaltade} & \nk{hjaltaddet}\\
        Past  & \nk{hjaltagne} & \nk{hjaltagnär} & \nk{hjaltagne} & \nk{hjaltagne}\\
    \end{tblr}
\end{table}

Adjectives generally appear before the noun they modify. Unlike in English, adjectives may be used to directly modify a personal or demonstrative pronoun or the name of a person. In such usage, the adjective must continue to agree with the personal or demonstrative pronoun or the name of the person in gender and number. (In English, this is accomplished by a relative clause although forms like `the happy I' are also possible. A definite article is not used in this construction in Niklorsk.)

\pex
\begingl
\gla Pisner hon stettatt sla kvett.//
\glb miserable.\F{} \Third{}\Sg{} cry.\Prog{} from yesterday//
\glft `The miserable she has been crying since yesterday.'//
\endgl
\xe

\subsection{Nominalization of adjectives}

The neuter form of an adjective can be used as a noun to refer to a person or thing that has the quality described by the adjective. If the adjective, however, refers back to a specific noun or pronoun, the adjective must agree in gender and number to the implicit or omitted noun.

\pex
\begingl
\gla Skune væd så. Skuner vædar så Matti.//
\glb blind see not blind.\F{} see.\Pst{} not Matti//
\glft `The blind do not see. The blind girl did not see Matti.'//
\endgl
\xe


\subsection{Comparison of adjectives}

The comparative form of an adjective is generally formed by adding \nk{-inn} or \nk{-ginn} to the base form of an adjective, according to whether the adjective ends in a consonant or a vowel, respectively. The comparative form is declined for gender and number, following the same pattern as the base form. The comparative ending triggers umlaut of the stem vowel in native words: written \nk{a}, \nk{o}, and \nk{u} alternate with \nk{ä}, \nk{ö}, and \nk{y}, respectively. These are stem alternations rather than a change to the ordinary pronunciation of the vowel letters.

The comparative suffix is stress-neutral. Umlaut changes the stem vowel without changing which syllable bears stress; \nk{bang} and \nk{bänginn}, for example, both have initial stress. As a productive Class W suffix, \nk{-inn/-ginn} preserves the base's accent category, subject to rhyme licensing (Chapter~\ref{chap:derivations}, \S\ref{sec:derivational-accent}).

The unstressed vowel of \nk{-inn/-ginn} undergoes weak-vowel reduction, while the umlaut of the stressed stem vowel remains distinct. Agreement endings on comparative adjectives follow the same reduction rules as those on positive adjectives: feminine \nk{-er} and plural \nk{-a} have [\ipa{ɐ}], and neuter \nk{-e} is assimilated or lost according to the preceding consonant. The comparative retains the consonantal material of its suffix.

\begin{table}[!ht]
    \caption{Comparative form of the adjective \nikp{bang}{good}}
    \label{tab:adj-comp}
    \small
    \begin{tblr}{
        colspec={XXXXX}
    }
        Base form & Masculine & Feminine & Neuter & Plural \\
        \nikp{bang}{good} & \nk{bänginn} & \nk{bänginner} & \nk{bänginne} & \nk{bänginna}\\
        
        \nikp{hot}{big} & \nk{hötinn} & \nk{hötinner} & \nk{hötinne} & \nk{hötinna}\\

        \nikp{kurt}{short} & \nk{kyrtinn} & \nk{kyrtinner} & \nk{kyrtinne} & \nk{kyrtinna}\\

    \end{tblr}
\end{table}

The point of comparison is introduced by \nikp{ins}{than.} The comparative form is used where English would use the base form to indicate that the quality described by the adjective is exhibited in the same degree by both referents (i.e., in English `as... as...' constructions). Here however, the adjective must be modified by the adverb \nikp{blä}{same, equally.}
\pex
\begingl
\gla Matti ö slössinn ins Andärs.//
\glb Matti be.\Pres{} tall.\Comp{} than Andärs//
\glft `Matti is taller than Andärs.'//
\endgl
\xe

\pex
\begingl
\gla Andärst stor ö blä slössinn ins Matti.//
\glb Andärs.\Gen{} brother be.\Pres{} same tall.\Comp{} than Matti//
\glft `Andärs' brother is as tall as Matti.'//
\endgl
\xe

\subsection{Possessive pronouns and similar forms}

The possessive forms of pronouns behave like adjectives and must agree with the noun they modify in gender and number. The possessive pronouns have slightly more irregular declensions than regular adjectives, as shown in Table~\ref{tab:possessive-declension}.

Their unstressed agreement endings undergo the same reduction, assimilation, and apocope as those of other adjectives. The irregular stem forms in the possessive paradigm remain distinct from these phonetic reductions.

\begin{table}[!ht]
    \caption{Declension of possessive pronouns}
    \label{tab:possessive-declension}
    \small
    \begin{tblr}{
        colspec={XXXXX}
    }
    & Masculine & Feminine & Neuter & Plural \\
    \nikp{ag}{I} & \nk{æga} & \nk{ægar} & \nk{æge} & \nk{æga}\\
    \nikp{meg}{you} & \nk{meg} & \nk{mör} & \nk{migge} & \nk{mega}\\
    \nikp{nar}{he} & \nk{næ} & \nk{nar} & \nk{næ} & \nk{na}\\
    \nikp{jö}{it} & \nk{jör} & \nk{jöre} & \nk{jö} & \nk{jör}\\
    \nikp{ga}{we} & \nk{saugna} & \nk{saugnar} & \nk{saugne} & \nk{saugna}\\
    \nikp{it}{you} & \nk{isse} & \nk{isser} & \nk{isse} & \nk{issa}\\
    \nikp{nör}{they} & \nk{nör} & \nk{nörer} & \nk{nöre} & \nk{nöra}\\
    \end{tblr}
\end{table}

\medskip

\noindent\begin{tabularx}{\linewidth}{XX}
    \nk{Saugna Tadd, käll dänn bærg\dots} & `Our Father, who art in heaven\dots'\\
\end{tabularx}

\section{Adverbs}

\subsection{Formation}

Adverbs are used to modify verbs, adjectives, or other adverbs. Most adverbs are formed by adding the suffix \nk{-æg/-gæg} to the base form of an adjective depending on whether the adjective ends in a consonant or a vowel, respectively. The adverbial form is invariable, except for an additional comparative form, which is formed by adding the suffix \nk{-inn} to the basic adverbial form. The comparative form is used to indicate that the action described by the verb is done to a greater extent than in comparison to some reference.

Both the adverbial suffix \nk{-æg/-gæg} and the comparative suffix \nk{-inn} preserve the stress of the base. Thus \nk{ront}, \nk{rontæg}, and \nk{rontæginn} all have initial stress (see Chapter~\ref{chap:derivations}, \S\ref{sec:affix-stress}). Productive formations preserve the base's accent category under the Class W rule in \S\ref{sec:derivational-accent}; an audible contrast still requires an eligible primary-stressed rhyme.

\pex
\a \nikp{maus}{happy} $\rightarrow$ \nikp{mausæg}{happily} $\rightarrow$ \nikp{mausæginn}{more happily}
\a \nikp{ront}{fast} $\rightarrow$ \nikp{rontæg}{quickly} $\rightarrow$ \nikp{rontæginn}{more quickly}
\a \nikp{tvent}{slow} $\rightarrow$ \nikp{tventæg}{slowly} $\rightarrow$ \nikp{tventæginn}{more slowly}
\a \nikp{gans}{rare} $\rightarrow$ \nikp{gansæg}{rarely} $\rightarrow$ \nikp{gansæginn}{more rarely}
\xe

Adverbs may be formed from present and past participles by adding the suffix \nk{-æg} to the participle form (as with regular adjectives) but with the final \nk{-e} of the participle dropped first.

\pex
\a \nikp{krugade}{befitting} $\rightarrow$ \nikp{krugadæg}{befittingly}, not \nk{krugadegæg}
\xe

Latinate adjectives ending in \nk{-al} form adverbs by taking the suffix \nk{-s}.

\begin{itemize}
    \item \nikp{national}{national} $\rightarrow$ \nikp{nationals}{nationally}
    \item \nikp{normal}{normal} $\rightarrow$ \nikp{normals}{normally}
    \item \nikp{personal}{personal} $\rightarrow$ \nikp{personals}{personally}
    \item \nikp{total}{total} $\rightarrow$ \nikp{totals}{totally}
    \item \nikp{optional}{optional} $\rightarrow$ \nikp{optionals}{optionally}
\end{itemize}

In neutral sentences, the adverb is placed directly after the verb it modifies. In an inverted sentence, the adverb is placed right before the verb.

\pex
\begingl
\gla Andärs naudar tventæg atta brord.//
\glb Andärs enter.\Pst{} slowly to:the room//
\glft `Andärs entered the room slowly.'//
\endgl
\xe

\pex
\begingl
\gla Tventæg naudar Andärs atta brord.//
\glb slowly enter.\Pst{} Andärs to:the room//
\glft `Slowly, Andärs entered the room.'//
\endgl
\xe

Some common adverbs are not derived from adjectives. They are listed below together with their comparative forms, most of them irregular. Some adverbs lack a comparative form altogether.

\pex
\a \nikp{bing}{well} $\rightarrow$ \nikp{bænginn}{better}
\a \nikp{augig}{often} $\rightarrow$ \nikp{auginn}{more often}
\a \nikp{när}{for a long time} $\rightarrow$ \nikp{närann}{longer}
\xe

\pex
Without a comparative form: \nikp{nörs}{never}, \nikp{asse}{always}, \nikp{lære}{now}, \nikp{nigot}{almost}, \nikp{gaups}{right now}, \nikp{bær}{somewhat}, \nikp{nosk}{very}, \nikp{bått}{very much}, \nikp{likke}{a little}, \nikp{ngå}{only}
\xe

\section{Prepositions}
\label{sec:prepositions}

\subsection{In general}

Prepositions introduce prepositional phrases and precede the complements they govern. A prepositional phrase may function as an adjunct, expressing such circumstances as time, place, or manner, or as an oblique argument licensed by a particular predicate.

The same phonological form may serve as a preposition in one construction and as a mobile verbal particle in another, as with \nikp{ån}{off, away}. These uses differ syntactically: a preposition heads a phrase containing its complement, whereas a mobile verbal particle forms part of a lexical predicate and obeys the particle-placement rules.

The relative positions of prepositional phrases, objects, and other clause constituents are described in Chapter~\ref{chap:syntax}, \S\ref{sec:pp-placement}. The distinction between mobile verbal particles and inseparable verbal prefixes is described in Chapter~\ref{chap:verbs}, \S\ref{sec:separable-verbs}.

\subsection{Location: \nk{i} and \nk{dänn}}
\label{sec:spatial-prepositions}

The general locative preposition \nikp{i}{in, on, at} identifies the geographical, temporal, circumstantial, or conceptual domain within which an entity or event is situated. It does not independently specify a bounded interior. Thus \nk{i Malmö} gives a geographical setting, while \nk{i kvett} gives an ordinary temporal setting. In conventional constructions with institutional destinations and activities, \nk{i} can also occur with motion predicates.

The delimited locative preposition \nikp{dänn}{within, inside} denotes inclusion within the bounds of a physical, institutional, perceptual, or eventive domain. The boundary need not be a wall: it can be the limits of a campus, the recognised premises of a cinema, the visual domain of a mirror, or the duration of a particular event. With ordinary places, \nk{dänn} generally corresponds to `inside' or `within'; with institutions and activities, it can express presence within their recognised setting without requiring literal physical enclosure.

\begin{table}[!ht]
    \caption{Delimited domains with \nk{dänn}}
    \label{tab:delimited-locations}
    \small
    \begin{tabularx}{\textwidth}{l X}
        \toprule
        \rowcolor{grammarTableHead}
        \textbf{Expression} & \textbf{Domain} \\
        \midrule
        \nk{dänn brug} & The physical enclosure of a house \\
        \nk{dänn brord} & The interior of a room \\
        \nk{dänn magasin} & The particular premises of a shop \\
        \nk{dänn kampus} & Conventionally bounded institutional grounds \\
        \nk{dänn kinnotek} & A cinema as a bounded venue \\
        \nk{dänn kung} & The interior or visual domain of a mirror \\
        \nk{dänn krust} & Participation within an earthquake event \\
        \bottomrule
    \end{tabularx}
\end{table}

Temporal uses preserve the distinction between general location and inclusion within a delimited domain. The preposition \nk{i} locates an event within a calendar period or ordinary temporal setting. The preposition \nk{dänn} presents an entity or event as participating within a bounded episode, event, or interval. Thus \nk{dänn krust i kvett} situates an event within an earthquake and places it on the previous day. A temporal interpretation of allative \nk{at} identifies a goal or endpoint rather than an interval containing the event (\S\ref{sec:goal-preposition}).

\subsection{Vertical position: \nk{være} and \nk{ud}}
\label{sec:vertical-prepositions}

The preposition \nikp{være}{under, below} locates an entity lower than the reference entity named by its complement. The preposition \nikp{ud}{over, above} expresses the corresponding higher position. Both precede their nominal complements. Locative predicates use \nk{vor}.

\pex
\begingl
\gla En glus vor være de kapp.//
\glb the.\M{} cat be.\Pres{} under the.\F{} coat//
\glft `The cat is under the coat.'//
\endgl
\xe

\pex
\begingl
\gla De skäpp vor ud eg skåv.//
\glb the.\F{} ship be.\Pres{} above the.\Nt{} reef//
\glft `The ship is above the reef.'//
\endgl
\xe

These prepositional phrases can also function as locative adjuncts. In a neutral transitive clause they follow the direct object; when fronted, they trigger the ordinary verb--subject inversion (Chapter~\ref{chap:syntax}, \S\ref{sec:pp-placement}).

\pex
\begingl
\gla Ag vædar de skäll være de skäpp.//
\glb \First{}\Sg{} see.\Pst{} the.\F{} fish below the.\F{} ship//
\glft `I saw the fish below the ship.'//
\endgl
\xe

\pex
\begingl
\gla Ud eg skåv vædar ag de skäll.//
\glb over the.\Nt{} reef see.\Pst{} \First{}\Sg{} the.\F{} fish//
\glft `Over the reef, I saw the fish.'//
\endgl
\xe

\subsection{Goals and endpoints: \nk{at}}
\label{sec:goal-preposition}

The general allative preposition \nikp{at}{to, into} introduces the goal or endpoint of directed motion, transfer, or another goal-oriented relationship. Its schematic semantic contribution is \(\operatorname{GOAL}(x,y)\), where \(x\) is the moving entity or transferred theme and \(y\) is the goal towards which the event is directed. The relationship need not involve literal physical motion.

\begin{table}[!ht]
    \caption{Goals introduced by \nk{at}}
    \label{tab:allative-goals}
    \small
    \begin{tabularx}{\textwidth}{X X}
        \toprule
        \rowcolor{grammarTableHead}
        \textbf{Construction} & \textbf{Goal and interpretation} \\
        \midrule
        \nk{malms at Norge} & Geographical destination: `to go to Norway' \\
        \nk{naudar atta brord} & Bounded destination: `entered the room' \\
        \nk{Ag brar afte at Matti.} & Recipient: `I gave a book to Matti.' \\
        \bottomrule
    \end{tabularx}
\end{table}

The translations `to' and `into' depend on the predicate and context. In \nk{naudar atta brord}, the verb \nikp{naud}{enter} specifies boundary crossing; \nk{at} identifies the destination, and the contracted article marks the noun phrase as definite. The form \nk{atta} has no independent ingressive meaning. In \nk{malms at Norge}, Norway is the destination; whether arrival is entailed depends on the aspectual interpretation of the whole predicate.

The contrast between \nk{at fjeltt} and \nk{atta fjeltt} concerns the presence of an overt definite article. The first has a bare nominal complement, while the second contains the article contracted with \nk{at}. A bare complement can still have a contextually identifiable referent, as in \nikp{at fjeltt}{to the party}. Article contraction and article omission are described in Chapter~\ref{chap:nouns}, \S\ref{sec:article-contraction}.

\subsection{Strong and weak forms}
\label{sec:preposition-reduction}

Prepositions precede their complements. Under contrastive focus they have strong forms with full vowel qualities. In ordinary unstressed speech they form a prosodic group with a neighbouring word and take the weak forms in Table~\ref{tab:weak-prepositions}. Their spelling and grammatical function remain unchanged by this attachment.

Weak prepositions have no independent lexical accent. A preposition pronounced as an independent stressed word follows the same rhyme-licensing conditions as other strong grammatical words (Chapter~\ref{chap:phonology}, \S\ref{sec:lexical-accent}).

\begin{table}[!ht]
    \caption{Weak forms of prepositions}
    \label{tab:weak-prepositions}
    \small
    \begin{tabularx}{\textwidth}{l X l X}
        \toprule
        \rowcolor{grammarTableHead}
        \textbf{Preposition} & \textbf{Meaning} & \textbf{Weak form} & \textbf{Rapid connected speech} \\
        \midrule
        \nk{at} & to & [\ipa{ət}] & [\ipa{t}] after a vowel \\
        \nk{i} & in, on, at & [\ipa{ɪ}] & [\ipa{j}] before a vowel \\
        \nk{av} & with & [\ipa{ɐʋ}] & [\ipa{ɐʋ}] \\
        \nk{sla} & from, of & [\ipa{slɐ}] & [\ipa{slɐ}] \\
        \nk{på} & by & [\ipa{pə}] & [\ipa{pə}] \\
        \nk{vin} & for, because of & [\ipa{vən}] & [\ipa{vn̩}] \\
        \bottomrule
    \end{tabularx}
\end{table}

The consonantal form of \nk{at} retains an audible /t/ and occurs where the connected-speech consonant sequence is licensed. Thus reduction does not remove recipient marking in \nikp{Ag brar afte at Matti}{I gave a book to Matti}; \nk{at} is required when the recipient follows the theme (see \S\ref{sec:ditransitive-verbs}). The agentive prepositions \nk{på} and \nk{par} remain distinct; the weak form [\ipa{pə}] belongs to \nk{på}. General weak-form processes and their relation to focus are described in Chapter~\ref{chap:phonology}, \S\ref{sec:connected-speech}.

\section{Conjunctions}

Conjunctions are used to connect words, phrases, or clauses. They can be divided into two main categories: coordinating conjunctions and subordinating conjunctions. Coordinating conjunctions link together words, phrases, or clauses of equal importance, while subordinating conjunctions are used to incorporate a dependent (or subordinate) clause into a main clause.

Below are some examples of coordinating conjunctions:

\pex
\a
\begingl
\gla Matti \emph{jä} Andärs bördar en bråst.//
\glb Matti and Andärs eat.\Pst{} the.\M{} cake//
\glft `Matti and Andärs ate the cake.'//
\endgl
\a
\begingl
\gla Ag malm at Tröndhem \emph{dårg} Matti malm at Rikke.//
\glb I go.\Pres{} to Trondheim but Matti go.\Pres{} to Rikke//
\glft `I am going to Trondheim but Matti is going to Rikke.'//
\endgl
\a
\begingl
\gla Meg igom komanders te \emph{tölle} kafe.//
\glb you.\Sg{} can.\Pres{} order.\Inf{} tea or coffee//
\glft `You can order tea or coffee.'//
\endgl
\xe

Similarly, some examples of subordinating conjunctions are given below:

\pex
\a
\begingl
\gla {Ny} slötar ag at färg, \ae{}ga stor svoratt någ.//
\glb when come.\Pst{} I to home, my.\M{} brother sleep.\Prog{} already//
\glft `When I came home, my brother was already asleep.'//
\endgl
\a
\xe


\subsection{Coordinating conjunctions}

Coordinating conjunctions generally appear between the words, phrases, or clauses they connect. The use of the conjunction does not usually affect the word order in the clause. The most common coordinating conjunctions are \nikp{jä}{and}, \nikp{dårg}{but} and \nikp{tölle}{or}.

\pex
\a
\begingl
\gla Martin jä Andärs slötäs at fjeltt.//
\glb Martin and Andärs come.\Fut{} to party//
\glft `Martin and Andärs will come to the party.'//
\endgl
\a
\begingl
\gla Ag bram vin exam dårg Matti pöljatt \ae{}gn.//
\glb I study.\Pres{} for exam but Matti disturb.\Prog{} me//
\glft `I am studying for the exam but Matti keeps on disturbing me.'//
\endgl
\xe

The adverb \nikp{sig}{as well, also} can be used with \nikp{jä}{and} to form a construction equivalent to `and also' in English. Similarly, \nikp{sig} can also be used with \nikp{dårg}, with the resulting construction equivalent to English `not only... but also'. Both \nikp{jä sig} and \nikp{dårg sig} cannot coordinate two complete sentences but only words, phrases, or clauses.

\pex
\begingl
\gla Ag jä sig Matti bram vin exam.//
\glb I and also Matti study.\Pres{} for exam//
\glft `I as well as Matti are studying for the exam.'//
\endgl
\xe

The construction \nk{vatt \dots{} jä (sig) \dots} corresponds to English `both \dots{} and \dots' or `both \dots{} as well as \dots'. Unlike English, however, this construction can coordinate more than two elements.

\pex
\begingl
\gla Vatt ag jä Martin malmar atta fjeltt.//
\glb both I and Martin go.\Pst{} to:the party//
\glft `Both I and Martin went to the party.'//
\endgl
\xe
\pex
\begingl
\gla Vatt Matti, Andärs jä sig Martin malmar atta fjeltt.//
\glb both Matti, Andärs and also Martin go.\Pst{} to:the party//
\glft `Matti, Andärs, and Martin all went to the party.'//
\endgl
\xe




